\lecture{4}
\title{Basic Kinematics}

\begin{document}

\subsection{Task vs. Configuration Space}

Recall the distinction between
\textbf{task space} and \textbf{configuration space}.
For example, given a Kuka IIWA manipulator\dots
\begin{itemize}
    \item The \textbf{task space} consists of
    just the position, orientation, and velocity of its end effector.
    \item The \textbf{configuration space}
    is much larger---it consists of all the
    angles and velocities at all $7$ of its joints.
\end{itemize}
\begin{center}
    \frame{\includegraphics[width=0.2\textwidth]{lecture-04/kuka-iiwa.jpeg}}
\end{center}
Last lecture, we discussed how we represent and work with
the \emph{task space} of the end effector.
In this lecture, we work in the \emph{configuration space}---how
should we move the joints such that the gripper
moves to the state we desire?

\textbf{Definition.}
We often talk about the following forms of kinematics.
\begin{itemize}
    \item \textbf{Forward Kinematics:}
    Given the joint positions, what is the end effector pose?
    \item \textbf{Inverse Kinematics:}
    Given the final pose, what are the joint positions?
    \item \textbf{Differential Kinematics:}
    Given the joint positions and velocities,
    what is the end effector pose and velocity?
    \item \textbf{Differential Inverse Kinematics:}
    Given the end effector pose and velocity,
    what are the joint positions and velocities?
\end{itemize}

\subsection{Forward Kinematics}

As one might expect, forward kinematics is
purely deterministic and very easy to compute:
just compose transforms.
To demonstrate, here's how that looks for
a simplified 2D Kuka IIWA.
\begin{center}
    \frame{\includegraphics[width=0.3\textwidth]{lecture-04/forward-kinematics.png}}
\end{center}
\[
    {}^aX^2(q_0, q_1)
    = {}^aX^0 \ {}^0X^b(q_0) \ {}^bX^1 \ {}^1X^c(q_1) \ {}^cX^2
\]
Some observations about the above:
\begin{itemize}
    \item The transforms
    $\{{}^aX^0, {}^bX^1, {}^cX^2\}$ are pure translations.
    \item The transforms
    $\{{}^0X^b(q_0), {}^1X^c(q_1)\}$
    are rotations parametrized by real numbers $(q_0, q_1)$.
    \item The superscripts on the RHS
    reads out a path $a \to 0 \to b \to 1 \to c \to 2$,
    which cleanly collapses to the LHS ${}^aX^2$.
\end{itemize}

\begin{remark}
    The pose of a single free body
    has six parameters---three for its position,
    and three for its orientation.
    But in Drake, the pose is represented
    using seven parameters, because the orientation
    is expressed using \emph{quaternions}.
\end{remark}

\subsection{Inverse Kinematics}

As one might expect, inverse kinematics
can have many solutions (or none) and is very hard to compute exactly.

One could use \emph{analytic} approaches
to find one of the (usually many!) \emph{exact} solutions
to an inverse kinematics problem.
But for systems with many parameters---such as
our Kuka IIWA manipulator---this is not a feasible approach.

Next lecture, we'll discuss \emph{numerical} approaches
for \emph{estimating} a solution to an inverse kinematics problem.

\subsection{Monogram Velocities}

Before moving on to differential kinematics,
we must discuss how we represent
velocity in monogram notation.

\textbf{Definition.}
The \textbf{spatial velocity} of $B$ relative to $A$ is:
\[ \frac{\mathrm{d}}{\mathrm{d}t} {}^AX^B_C
= {}^AV^B_C = \begin{bmatrix}
{}^A\omega^B_C \\ {}^Av^B_C
\end{bmatrix}
= \begin{bmatrix}
{}^A\omega^B_{C, x} \\
{}^A\omega^B_{C, y} \\
{}^A\omega^B_{C, z} \\
{}^Av^B_{C, x} \\
{}^Av^B_{C, y} \\
{}^Av^B_{C, z}
\end{bmatrix} \in \mathbb{R}^6, \]
where ${}^A\omega^B_C$ is angular velocity,
and ${}^Av^B_C$ is translational velocity.

\begin{remark}
    Note that angular velocity ${}^A\omega^B_C$
    is simply a three-number vector---as in 8.01,
    it points along the axis of rotation,
    with magnitude given by the speed.
    So ${}^AV^B_C$ is a six-number vector.
\end{remark}

\begin{theorem}{Monogram Velocities}
    The following identities are all true.
    \begin{itemize}
        \item \emph{(Change of Frame)} We have
        ${}^Av^B_G = {}^GR^F {}^Av^B_F$ and
        ${}^A\omega^B_G = {}^GR^F {}^A\omega^B_F$.
        \item \emph{(Additivity)} We have
        ${}^A\omega^B_F + {}^B\omega^C_F = {}^A\omega^C_F$ and
        ${}^Av^C_F = {}^Av^B_F + {}^Bv^C_F
        + {}^A\omega^B_F \times {}^Bp^C_F$.
        \item \emph{(Inverse)} We have
        ${}^A\omega^B_F = -{}^B\omega^A_F$ and
        ${}^Av^B_F = -{}^Bv^A_F -
        {}^A\omega^B_F \times {}^Bp^A_F$.
    \end{itemize}
\end{theorem}

\begin{proof}
    The simplified 2D Kuka IIWA example from earlier
    is a prototype to use when
    thinking through these identities.

    For the identity ${}^Av^C_F = {}^Av^B_F + {}^Bv^C_F
    + {}^A\omega^B_F \times {}^Bp^C_F$, for example,
    recall that ${}^Bv^C_F$ only captures
    the movement of joint $C$ in the frame of joint $B$.
    So if only joint $B$ is rotating, then ${}^Bv^C_F$ is zero,
    even though $C$ is moving in the frame of $A$.

    To account for the movement of joint $C$ caused by
    the rotation of frame $B$, we must add the term
    ${}^A\omega^B_F \times {}^Bp^C_F$.
    (You can check that the proportionalities make sense.)
    ~ $\blacksquare$
\end{proof}

\subsection{Differential Kinematics: Analytic Jacobian}

Back to something easy.
From forward kinematics, we know that
for any body $B$ with $k$ parameters $q$
that define its pose, there is a map
$f_{\mathrm{kin}}^B \colon \mathbb{R}^k \to \mathbb{R}^{3 + \ell}$
such that $X^B = f_{\mathrm{kin}}^B(q)$.
We can now \emph{differentiate} this map to achieve:
\[
\mathrm{d}X^B =
\frac{\partial f_{\mathrm{kin}}^B(q)}{\partial q} \mathrm{d}q
= J^B(q) \mathrm{d}q,
~ \text{ where } J^B(q) \text{ is the } \textbf{Analytic Jacobian.}
\]
In particular, the analytic Jacobian is
a derivative of a pose parametrization.  Explicitly,
\[ x(q) = \begin{bmatrix} p(q) \\ \phi(q) \end{bmatrix}
\text{ where } p(q) \in \mathbb{R}^3
\text{ and } \phi(q) \in \mathbb{R}^{\ell}, \]
where $\ell$ depends on the choice of orientation parametrization.
For example, roll-pitch-yaw would require $\ell = 3$,
whereas quaternions would yield $\ell = 4$,
and rotation matrices would yield $\ell = 9$.

It then follows that $J^B(q)$ and $\mathrm{d}q$ look like:
\[ J^B(q) = \frac{\partial x}{\partial q} = \begin{bmatrix}
\partial p/\partial q_1 & \partial p/\partial q_2
& \cdots & \partial p/\partial q_k \\
\partial \phi/\partial q_1 & \partial \phi/\partial q_2
& \cdots & \partial \phi/\partial q_k \\
\end{bmatrix}
\in \mathbb{R}^{(3 + \ell) \times (k)}
~ \text{ and } ~
\mathrm{d}q = \begin{bmatrix}
\mathrm{d}q_1 \\ \mathrm{d}q_2 \\ \vdots \\ \mathrm{d}q_k
\end{bmatrix} \in \mathbb{R}^k. \]

\subsection{Differential Kinematics: Geometric Jacobian}

There's a nice simplifying trick, though.
Whereas orientation has many different parametrizations---all
of them flawed---velocity (as in,
translational and rotational velocity combined)
has a very natural $6$-variable parametrization:
the spatial velocity.

\begin{remark}
    Recall that ${}^A\omega^B_C \in \mathbb{R}^3$
    points along the axis of rotation with magnitude given by
    the angular speed.  Equivalently, if
    \[ {}^A\omega^B_C = \begin{bmatrix} {}^A\omega^B_{C, x} \\
    {}^A\omega^B_{C, y} \\
    {}^A\omega^B_{C, z} \end{bmatrix} \in \mathbb{R}^3, \]
    that means the rotation is equivalent to the composition of
    a rotation with speed ${}^A\omega^B_{C, x}$ about the $x$-axis,
    a rotation with speed ${}^A\omega^B_{C, y}$ about the $y$-axis,
    and a rotation with speed ${}^A\omega^B_{C, z}$ about the $z$-axis.
\end{remark}

Then $J^B(q)$ is not very hard to compute;
the $i^{\text{th}}$ column is an $\mathbb{R}^6$-vector
that reads the rotational and translational velocities
incurred by changing the $i^{\text{th}}$ parameter $q_i$.

For example, for \textbf{revolute joints}
(i.e., every joint $q_i$ can only rotate about a line $\hat{z}_i$), we have:
\[
    J^B(q) = \begin{bmatrix}
    \hat{z}_1 & \hat{z}_2 & \cdots & \hat{z}_k \\
    \hat{z}_1 \times {}^{J_1}p^B & \hat{z}_2 \times {}^{J_2}p^B & \cdots &
    \hat{z}_k \times {}^{J_k}p^B
    \end{bmatrix} \in \mathbb{R}^{6 \times k}
    ~ \text{ and } ~
    \mathrm{d}q = \begin{bmatrix}
        \mathrm{d}q_1 \\ \mathrm{d}q_2 \\ \vdots \\ \mathrm{d}q_k
    \end{bmatrix} \in \mathbb{R}^k.
\]
When parametrizing velocities in this way,
the $J^B(q)$ we achieve is specifically called
the \textbf{Geometric Jacobian}.

Notice how each joint gets one column, and the
height of the geometric Jacobian is always $6$,
independent of how we internally represent
the end effector's orientation.

\begin{problem}
    Recall the example from earlier,
    specifically in the case where both rods have length $L = 1$.
    \begin{center}
        \frame{\includegraphics[width=0.3\textwidth]{lecture-04/forward-kinematics.png}}
    \end{center}
    What is the geometric Jacobian $J^C(q)$ of the tip $C$?
\end{problem}

\begin{solution}
    There are $k = 2$ joints, and they are both revolute joints.
    Applying the formula above, we find:
    \[
    J^C(q) = \begin{bmatrix}
    0 & 0 \\
    0 & 0 \\
    1 & 1 \\
    -\sin(q_0) - \sin(q_0 + q_1) & -\sin(q_0 + q_1) \\
    \cos(q_0) + \cos(q_0 + q_1) & \cos(q_0 + q_1) \\
    0 & 0
    \end{bmatrix} \in \mathbb{R}^{6 \times 2}.
    \]
    Check this for yourself.
    Note that the angular speed is not factored into $J^C(q)$,
    because it's factored into $\mathrm{d}q$ instead. ~ $\blacksquare$
\end{solution}

You can also check for yourself that
the dimensions of the geometric Jacobian of the IIWA are $6 \times 7$.

\subsection{Inverse Differential Kinematics}

Inverse differential kinematics is actually not so hard;
from the relationship $\mathrm{d}X^B = J^B(q) \mathrm{d}q$,
it follows that $\mathrm{d}q = J^B(q)^{-1}\mathrm{d}X^B$.

But this doesn't actually work as written,
because $J^B(q)^{-1}$ need not exist,
because the geometric Jacobian might not be square.
(In general, its dimensions are $6 \times k$,
where $k$ is the number of joints.)

The fix, in general, is to take inspiration from
\href{/18.650/lectures/23}{Lecture 23} of 18.650, on linear regression.
\begin{itemize}
    \item In linear regression,
    we seek parameters $\beta \in \mathbb{R}^k$ such that
    $\mathbb{X}\beta \in \mathbb{R}^n$ is as close to $Y \in \mathbb{R}^n$
    as the column space of $\mathbb{X}$ allows.
    \begin{itemize}
        \item As computed in 18.650, the desired parameter is
        $\beta = (\mathbb{X}^{\top}\mathbb{X})^{-1}\mathbb{X}^{\top}Y$.
    \end{itemize}
    \item In this class, we seek an input $\mathrm{d}q$ such that
    $J \cdot \mathrm{d}q$ is as close to $\mathrm{d} X^B$
    as the column space of $J$ allows.
    \begin{itemize}
        \item Expectedly, the desired input is
        $\mathrm{d}q = (J^{\top}J)^{-1}J^{\top} \mathrm{d} X^B$.
    \end{itemize}
\end{itemize}

\textbf{Definition.}
Given a matrix $J$ with full row rank, the \textbf{Moore-Penrose Pseudo-Inverse} of $J$
is given by $J^+ := J^{\top}(JJ^{\top})^{-1}$.

\begin{remark}
    When $J$ instead has full column rank (as in regression),
    $J^+ = (J^{\top}J)^{-1}J^{\top}$;
    the two formulas agree only when $J$ is square and invertible.
    For the $6 \times 7$ IIWA Jacobian, $J^{\top}J$ is never invertible.
\end{remark}

And here's the property of $J^+$ that makes it so useful,
as in both 18.650 and this class.

\begin{theorem}{Pseudo-Inverse Approximation}
    For any matrix $J$ and vector $b$ (with appropriate dimensions),
    the vector $x^* := J^+b$ is the vector with minimum magnitude
    across all vectors that minimize $\|Jx - b\|$.  In particular,
    \begin{itemize}
        \item If the column space of $J$ does not contain $b$,
        then $x^* := J^+b$ is such that $\|Jx - b\|$ is minimized.
        \item If the column space of $J$ contains $b$
        (that is, $\|Jx - b\|$ can be zero),
        then $x^* := J^+b$ is the solution to $Jx = b$ minimizing $\|x\|$.
    \end{itemize}
\end{theorem}

But there's still an issue with the pseudo-inverse
as stated above; its values can explode near \textbf{singularities}.

\textbf{Definition.}
We say there is a \textbf{singularity} when
the rows of $J^B(q)$ for the velocities we want to control
are not full row rank.

\textbf{Example.}
In the two-rod example from earlier,
when considering only the gripper's $(x, y)$ velocity,
one singularity occurs when $q_1 = 0$:
both joints can then only move the gripper perpendicular to the arm.

The fix for this problem is also inspired by 18.650;
use \textbf{damped pseudo-inverse} (or equivalently,
add a ridge regression term).
\[
    J^* = J^{\top}(JJ^{\top} + \lambda I)^{-1}.
\]
The $\lambda I$ term prevents things from exploding
when $JJ^{\top}$ is close to zero in some directions.

\begin{remark}
    The existence of a singularity
    is not enough to cause concern
    if that singularity occurs momentarily.

    For example, if $q_1$ gradually decreases
    from $+10^{\circ}$ to $-10^{\circ}$,
    then there is a moment when $q_1 = 0^{\circ}$
    and a singularity occurs.  But such a moment is instantaneous;
    the pre-existing momentum of the robot
    gets it to pass the singularity unharmed.
\end{remark}

\end{document}